\documentclass[11pt]{article}
\usepackage[a4paper,margin=25mm]{geometry}
\usepackage[T1]{fontenc}
\usepackage{lmodern,microtype,amsmath,amssymb,amsthm,mathtools}
\usepackage{xcolor}
\usepackage[colorlinks=true,linkcolor=blue!50!black,citecolor=blue!50!black,urlcolor=blue!50!black]{hyperref}
\setlength{\emergencystretch}{2em}
\newtheorem{theorem}{Theorem}[section]
\newtheorem{proposition}[theorem]{Proposition}
\newtheorem{corollary}[theorem]{Corollary}
\newtheorem{lemma}[theorem]{Lemma}
\theoremstyle{remark}\newtheorem{remark}[theorem]{Remark}
\newcommand{\Ric}{\operatorname{Ric}}
\newcommand{\tr}{\operatorname{tr}}
\newcommand{\diver}{\operatorname{div}}
\newcommand{\Sym}{\operatorname{Sym}}
\newcommand{\Id}{\operatorname{Id}}
\newcommand{\dd}{\,dV_g}
\title{Q4: Torsion Zeros and Transverse Weyl Curvature\\[3pt]
\large Local structure, positive transverse indices, and the remaining global issue}
\author{Research continuation from the supplied four-dimensional Einstein manuscripts}
\date{9 October 2026}
\begin{document}
\maketitle
\begin{abstract}
For an almost-K\"ahler Einstein four-manifold we derive an exact
second-derivative identity at every torsion zero. Combining this identity
with the adapted-frame curvature equations gives a complex normal form
for the derivative of the torsion, a nonnegative real determinant, and a
sharp algebraic estimate involving the transverse Weyl component. In
particular, a torsion zero where that component is nonzero is isolated and
transverse, with positive index. If all zeros are transverse on a closed
manifold, their number is $5\chi+6\tau$. These are necessary conditions for
an actual solution, not an existence theorem for prescribed jets. They do
not establish the closed negative Goldberg conclusion.
\end{abstract}

\section{Conventions and the precise result}
Let $(M^4,g,J,\Omega)$ be a smooth, positive-definite, oriented
almost-K\"ahler Einstein manifold. Thus
\[
 d\Omega=0,\qquad \Omega(X,Y)=g(JX,Y),\qquad |\Omega|^2=2,
 \qquad \Ric=\frac{s}{4}g.
\]
The scalar $s$ is constant. The local results do not require completeness,
compactness, or a sign for $s$. Write $W=W^+$ as a symmetric trace-free
endomorphism of $\Lambda^+$, with Hilbert--Schmidt operator norm. Set
\begin{equation}
 q=|\nabla\Omega|^2,\quad n=\Omega/\sqrt2,\quad
 t=\langle Wn,n\rangle=\frac{s}{6}+\frac q4,\quad
 \eta=W\Omega-t\Omega,\quad v=\eta/\sqrt2.
 \label{eq:defs}
\end{equation}
With respect to $\mathbb Rn\oplus n^\perp$,
\begin{equation}
 W=\begin{pmatrix}t&v^*\\v&-\frac t2\Id+V\end{pmatrix},
 \qquad \tr V=0,\qquad
 |W|^2=\frac32t^2+|\eta|^2+|V|^2.
 \label{eq:Wblock}
\end{equation}
Here $V$ is the $W_3^+$ component in the chosen almost-Hermitian
splitting, not the full restriction $W|_{n^\perp}$.

Locally choose $\phi,J\phi$ of norm $\sqrt2$, orthogonal to $\Omega$.
The adapted-frame equations of \cite[Section 2.3]{AAD} are
\begin{align}
 \nabla\Omega&=a\otimes\phi-Ja\otimes J\phi,& q&=4|a|^2,\label{eq:torsion}\\
 da-Ja\wedge b&=-\mathcal R(J\phi),&
 d(Ja)+a\wedge b&=-\mathcal R(\phi).\label{eq:frame}
\end{align}
The covector convention is $(Ja)(X)=-a(JX)$.
At a zero $p$ of $a$, let
\[
 L:T_pM\longrightarrow T_pM,\qquad L(X)=(\nabla_Xa)^\sharp.
\]
This derivative is independent of derivatives of the chosen frame at $p$.
Changing $\phi$ multiplies $L$ on its target by a unitary rotation, so its
real determinant, singular values, and vanishing rank are well defined.

\begin{theorem}[Local structure of a torsion zero]\label{thm:main}
At every point $p$ where $q(p)=0$:
\begin{equation}
 \boxed{|\nabla^2\Omega|^2=4|V|^2-6|\eta|^2.}\label{eq:zero-balance}
\end{equation}
There are complex orthonormal identifications of the source and target
of $L$ with $\mathbb C^2$ in which
\begin{equation}
 L(z)=c z+A\bar z,\qquad c\in\mathbb C,\quad A\in M_2(\mathbb C),
 \qquad \tr(A\bar A)=-8|c|^2.\label{eq:normal}
\end{equation}
The bar in $A\bar A$ is entrywise conjugation, not the Hermitian adjoint.
Furthermore,
\begin{align}
 |\eta|^2&=8|c|^2,\label{eq:eta-c}\\
 \boxed{|V|^2-4|\eta|^2=4\|\Sym A\|_F^2\ge0,}\label{eq:Veta}\\
 \boxed{\det_{\mathbb R}L
 =9|c|^4+|\det_{\mathbb C}A|^2
 \ge25|c|^4=\frac{25}{64}|\eta|^4.}\label{eq:det}
\end{align}
Consequently, if $\eta(p)\ne0$, the torsion zero is isolated and
transverse, and $q$ has a nondegenerate local minimum there. Every
transverse torsion zero has positive local index. If $L$ is singular,
then $\eta(p)=0$ and $\operatorname{rank}_{\mathbb R}L\in\{0,2\}$.
\end{theorem}

\section{The second-derivative identity}
Use the nonnegative scalar Laplacian $\Delta=-\diver\nabla$ and the
nonnegative rough Laplacian $\nabla^*\nabla$. Since $\Omega$ is a
self-dual harmonic form,
\begin{equation}
 \nabla^*\nabla\Omega=2W\Omega-\frac{s}{3}\Omega
                    =\frac q2\Omega+2\eta.\label{eq:form-WB}
\end{equation}
This also proves the expression for $t$ in \eqref{eq:defs}. Einstein
metrics satisfy $\delta W=0$, and the classical Weyl--Weitzenb\"ock
identity is
\begin{equation}
 \nabla^*\nabla W+\frac{s}{2}W-6W^2+2|W|^2\Id=0.
 \label{eq:W-WB}
\end{equation}
See \cite{Derdzinski}; the same operator normalization appears in
\cite[Lemma 2.5, with $f=1$]{CD}.

At a torsion zero, $\nabla n=0$, $t=s/6$, and
$\nabla^*\nabla n=2v$. Differentiating $t=\langle Wn,n\rangle$
therefore gives
\begin{align*}
 \Delta t
 &=\langle(\nabla^*\nabla W)n,n\rangle
          +2\langle W(\nabla^*\nabla n),n\rangle\\
 &=6(t^2+|v|^2)-\frac{s}{2}t-2|W|^2+4|v|^2\\
 &=3t^2-\frac{s}{2}t+6|v|^2-2|V|^2\\
 &=3|\eta|^2-2|V|^2.
\end{align*}
Since $\Delta q=4\Delta t$, whereas
$\Delta|\nabla\Omega|^2=-2|\nabla^2\Omega|^2$ at a zero of
$\nabla\Omega$, this proves \eqref{eq:zero-balance}. All derivatives
here are full covariant derivatives; $\nabla^2\Omega$ is not merely
its skew-symmetrization.

\section{The complete first derivative and its determinant}
Choose an oriented orthonormal coframe at $p$ with
\[
 \Omega=e^{12}+e^{34},\quad
 \phi=e^{13}-e^{24},\quad J\phi=e^{14}+e^{23}.
\]
At $a(p)=0$, also $\nabla J(p)=0$. Einstein curvature preserves
$\Lambda^+$, so \eqref{eq:frame} shows that $da$ and $d(Ja)$ are
self-dual at $p$. These give six independent real linear equations on
$\nabla a$. Equating their anti-self-dual components gives the following
complete ten-parameter matrix. Put
\[
 c=\ell+i m,\qquad
 A=\begin{pmatrix}p+i q_0&r+i t_0\\u+i v_0&w+i z_0\end{pmatrix}.
\]
The matrix with derivative index on rows, namely $L^{\mathsf T}$, is
\begin{equation}
 \begin{pmatrix}
 \ell+p&m+q_0&u&v_0\\
 -m+q_0&\ell-p&v_0&-u\\
 r&t_0&\ell+w&m+z_0\\
 t_0&-r&-m+z_0&\ell-w
 \end{pmatrix}.\label{eq:real-matrix}
\end{equation}
Its transpose is exactly the real matrix of $z\mapsto cz+A\bar z$.
Thus no additional modes have been discarded.

Write $\mathcal R^+=W+s\Id/12$. At the zero, the transverse
$2\times2$ block of $\mathcal R^+$ is $V$, because $t=s/6$.
From $\mathcal R(\phi)=-d(Ja)$ and
$\mathcal R(J\phi)=-da$ one obtains
\[
 (\mathcal R^+)_{12}=-2\ell,\qquad
 (\mathcal R^+)_{13}=-2m,\qquad
 V=\begin{pmatrix}v_0-t_0&r-u\\r-u&t_0-v_0\end{pmatrix}.
\]
Hence
\begin{equation}
 |\eta|^2=8|c|^2,\qquad
 |V|^2=2|A_{12}-A_{21}|^2,\qquad
 |L|^2=4|c|^2+2\|A\|_F^2.\label{eq:norms}
\end{equation}
Differentiating \eqref{eq:torsion} at $a=0$ gives
$|\nabla^2\Omega|^2=4|L|^2$. Thus \eqref{eq:zero-balance}
is precisely
\[
 8|c|^2+\|A\|_F^2-|A_{12}-A_{21}|^2=0.
\]
The last two terms equal $\tr(A\bar A)$, which proves
\eqref{eq:normal}.

Let $A=S+K$, where $S$ is complex symmetric and $K$ is complex skew-symmetric
(that is, $S=(A+A^{\mathsf T})/2$). Then
\[
 \tr(A\bar A)=\|S\|_F^2-\|K\|_F^2,
 \qquad |V|^2=4\|K\|_F^2.
\]
Equation \eqref{eq:normal} yields
$\|K\|_F^2=\|S\|_F^2+8|c|^2$, and proves \eqref{eq:Veta}.
In particular the norm identity also reads
\[
 |\nabla^2\Omega|^2=10|\eta|^2+16\|S\|_F^2.
\]

The real determinant of a complex-linear plus anti-linear map is
\begin{align*}
 \det_{\mathbb R}(c\Id+A\mathsf C)
 &=\det_{\mathbb C}
       \begin{pmatrix}c\Id&A\\\bar A&\bar c\Id\end{pmatrix}\\
 &=|c|^4-|c|^2\tr(A\bar A)+|\det A|^2,
\end{align*}
where $\mathsf C$ is complex conjugation. The equality holds for all
$c$, by polynomial continuation if necessary. Substitution proves the
first equality in \eqref{eq:det}.

For $A=\left(\begin{smallmatrix}a&b\\d&e\end{smallmatrix}\right)$,
\[
 \tr(A\bar A)=|a|^2+|e|^2+2\Re(b\bar d)
 \ge2|ae|-2|bd|\ge-2|ae-bd|.
\]
It follows that $|\det A|\ge4|c|^2$, which proves the lower bound in
\eqref{eq:det}. Equality in both \eqref{eq:Veta} and \eqref{eq:det}
is possible for the algebraic data
\[
 c\in\mathbb R,\qquad A=\begin{pmatrix}0&2c\\-2c&0\end{pmatrix}.
\]
This establishes sharpness of the algebraic coefficients under the
necessary equations above; it does not construct an Einstein metric
realizing these prescribed data.

If $\eta\ne0$, the determinant is positive. The inverse function theorem
applied to $a$ gives an isolated transverse zero; since $q=4|a|^2$,
\[
 \operatorname{Hess}_p q(X,Y)=8\langle LX,LY\rangle.
\]
If $\det L=0$, \eqref{eq:det} forces $c=0$ and $\det A=0$.
Then $L=A\mathsf C$ has real rank zero or two. This completes the local
algebraic assertions of Theorem~\ref{thm:main}.

\section{Indices on a closed manifold}
Use the standard Nijenhuis convention
\[
 N_J(X,Y)=[JX,JY]-J[JX,Y]-J[X,JY]-[X,Y].
\]
The adapted-frame expression for $\nabla J$ implies
\[
 N_J(e_1,e_3)=2J a^\sharp.
\]
After complexifying this is an orientation-preserving linear
identification of the four real torsion components with the complex
rank-two bundle
\[
 \mathcal N=\Lambda^{0,2}T^*M\otimes T^{1,0}M.
\]
Changing the adapted frame changes this identification by
orientation-preserving complex transformations. At a zero, derivatives
of the identification itself multiply $a=0$ and do not affect the
Jacobian. Thus the sign of a transverse local Nijenhuis index is the
sign of $\det L$. By \eqref{eq:det}, it is positive.

\begin{corollary}[Count when the torsion zeros are transverse]
If $M$ is closed and every zero of $N_J$ is transverse, then
\begin{equation}
 \boxed{\#\{N_J=0\}=5\chi(M)+6\tau(M).}\label{eq:count}
\end{equation}
In particular the same formula holds if $\eta$ is nonzero at every
torsion zero. If a torsion zero is degenerate, then $\eta$ vanishes
there and the first derivative has real rank zero or two.
\end{corollary}
\begin{proof}
The zeros are isolated and finite. If $x_1,x_2$ are the formal Chern
roots of $T^{1,0}M$, those of $\mathcal N$ are
$2x_1+x_2$ and $x_1+2x_2$. Therefore
\[
 c_2(\mathcal N)=2c_1^2(TM,J)+c_2(TM,J)=5\chi+6\tau.
\]
The index theorem for a transverse section equates this Chern number
to the sum of local indices. Each is $+1$ by Theorem~\ref{thm:main}.
The characteristic-class formula is the classical input used in the
project's torsion-zero argument, following \cite{Armstrong}.
\end{proof}

For a closed Einstein metric with $s\ne0$, Chern--Weil gives
$\chi>0$ and $2\chi+3\tau>0$, so
$5\chi+6\tau=\chi+2(2\chi+3\tau)>0$.
Consequently \eqref{eq:count} is compatible with Einstein topology.
It is not a contradiction. Transversality is an extra property of the
actual Nijenhuis section and is not obtained by perturbing that section
while discarding the almost-K\"ahler Einstein equations.

\section{A scalar current identity and its limit}
For clarity, the pointwise identity behind the usual integrated balance
can also be retained without dropping its current. Define the one-form
\[
 \theta(X)=\langle\eta,\nabla_X\Omega\rangle.
\]
Then the same Einstein and harmonic-form equations imply
\begin{equation}
 \boxed{\frac14\Delta q+2\diver(\theta^\sharp)
       +\frac{s}{8}q+\frac{3}{16}q^2+|\eta|^2+2|V|^2=0.}
 \label{eq:current}
\end{equation}
Here is a direct derivation. Let $T_i=\nabla_i n$, $f=\sum_i|T_i|^2=q/2$,
and $\beta_i=\langle v,T_i\rangle=\theta_i/2$.
The almost-K\"ahler form of \eqref{eq:torsion} gives
\[
 \sum_iT_i\otimes T_i=\frac f2(\Id-n\otimes n),\quad
 \sum_i\langle WT_i,T_i\rangle=-\frac{tf}{2},\quad
 \nabla^*\nabla n=fn+2v.
\]
Product differentiation using \eqref{eq:W-WB} first yields
\[
 \Delta t=3t^2-\frac{s}{2}t+6|v|^2-2|V|^2+3tf
            -4\sum_i\langle(\nabla_iW)T_i,n\rangle.
\]
A separate differentiation of $\beta$ gives
\[
 \sum_i\langle(\nabla_iW)T_i,n\rangle
       =\diver\beta+\frac32tf+2|v|^2.
\]
Since $f=2t-s/3$, their combination is
\[
 \Delta t=\frac{s}{2}t-3t^2-2|v|^2-2|V|^2-4\diver\beta.
\]
Substitution of \eqref{eq:defs} proves \eqref{eq:current}.

On a closed manifold, integration recovers the balance already contained
in the supplied manuscript:
\[
 \int\left(|\eta|^2+2|V|^2+\frac{3q^2}{16}\right)\dd
 =-\frac{s}{8}\int q\dd.
\]
That integrated identity is not an additional constraint introduced here.
At a torsion zero, $\theta=0$ but
$\diver\theta=-2|\eta|^2$, so dropping the divergence pointwise would be
incorrect. Keeping it recovers \eqref{eq:zero-balance}.

\paragraph{Connection to Q4.}
The supplied manuscript previously established that a closed
nonzero-scalar Einstein almost-K\"ahler candidate has a torsion zero.
The result here specifies the full first derivative at such a zero:
nonzero transverse Weyl curvature forces a positive isolated zero, whereas
a degenerate zero must lie in the locus $\eta=0$ and has even derivative
rank. None of these statements proves that $\eta$ vanishes on all of $M$.
The compact third-Gray theorem in \cite{AAD} would finish the integrability
argument if that global vanishing were obtained. A pointwise zero of the
torsion is not, by itself, the infinite-order vanishing required by a
strong unique-continuation argument.

\paragraph{Reproducible finite calculations.}
The accompanying \texttt{verify\_zero\_structure.py} checks the six linear
self-duality conditions, their ten-dimensional solution space, the norm
identities, the determinant formula, and the algebraic equality case.
It also substitutes the project's explicit rational noncompact metric
into \eqref{eq:current}; the rational residual is exactly zero. That is a
consistency check on a previously supplied model, not a new example or a
metric-existence argument for the algebraic data at a torsion zero.

\begin{thebibliography}{9}
\bibitem{AAD}
V. Apostolov, J. Armstrong and T. Dr\u aghici,
\emph{Local models and integrability of certain almost K\"ahler 4-manifolds},
Math. Ann. \textbf{323} (2002), 633--666.
\href{https://arxiv.org/abs/math/0109098}{arXiv:math/0109098}.
Equations (21)--(24), the harmonic-form identities, and Theorem 2.
\bibitem{Derdzinski}
A. Derdzi\'nski,
\emph{Self-dual K\"ahler manifolds and Einstein manifolds of dimension four},
Compositio Math. \textbf{49} (1983), 405--433.
\href{https://archive.numdam.org/item/CM_1983__49_3_405_0/}{Journal archive}.
\bibitem{CD}
G. Catino and D. Dameno,
\emph{Conformal K\"ahler Rigidity of Einstein Four-Manifolds},
2026 preprint, Lemma 2.5.
\href{https://arxiv.org/abs/2607.21538}{arXiv:2607.21538}.
Only the classical Weyl--Weitzenb\"ock formula, with $f=1$, is used here.
\bibitem{Armstrong}
J. Armstrong, \emph{Almost-K\"ahler Geometry}, doctoral thesis.
\href{https://johnarmstrongmaths.com/publications/resources/thesis.pdf}{Author's text}.
Characteristic class of the Nijenhuis bundle.
\bibitem{Project}
Supplied project manuscripts, 8 October 2026:
\texttt{01\_Research\_Synthesis(1).tex}, the Goldberg section;
\texttt{02\_Mathematical\_Source\_Companion(1).tex}, Step 55;
and \texttt{03\_Subsequent\_Mathematical\_Developments(1).tex},
the explicit Einstein example and the closed torsion-zero proposition.
These identify the existing project results to which the local calculation
is connected. No FD, N1, or global Einstein-gap theorem is used in this proof.
\end{thebibliography}
\end{document}
